\documentclass[11pt,reqno]{amsart}
\usepackage[T1]{fontenc}
\usepackage[utf8]{inputenc}
\usepackage{mathpazo}
\usepackage{amsmath,amssymb,amsthm,mathtools}
\usepackage{microtype}
\usepackage{geometry}
\usepackage{enumitem}
\usepackage{xcolor}
\usepackage{hyperref}
\geometry{a4paper,margin=30mm,headheight=15pt}
\hypersetup{colorlinks=true,linkcolor=black,citecolor=black,urlcolor=blue!45!black,pdfauthor={Arseniy Akopyan},pdftitle={A global isoperimetric inequality for spherical tetrahedra}}
\allowdisplaybreaks[1]
\setlength{\parskip}{0pt}
\setlength{\parindent}{15pt}
\newtheorem{theorem}{Theorem}[section]
\newtheorem{lemma}[theorem]{Lemma}
\newtheorem{proposition}[theorem]{Proposition}
\newtheorem{corollary}[theorem]{Corollary}
\newtheorem{claim}[theorem]{Claim}
\theoremstyle{remark}
\newtheorem{remark}[theorem]{Remark}
\DeclareMathOperator{\area}{area}
\DeclareMathOperator{\vol}{vol}
\newcommand{\Sph}{\mathbb S}
\newcommand{\R}{\mathbb R}
\newcommand{\dd}{\,\mathrm d}
\newcommand{\ip}[2]{\langle #1,#2\rangle}
\newcommand{\tr}{\operatorname{tr}}
\renewcommand{\arraystretch}{1.15}
\title[A global inequality for spherical tetrahedra]{A global isoperimetric inequality\\for spherical tetrahedra}
\author{Arseniy Akopyan}
\date{November 2026}
\subjclass[2020]{52A40, 52B10, 51M20, 53C20}
\keywords{Spherical tetrahedron, isoperimetric inequality, Schlaefli formula, Gram matrix, spherical polygon}
\begin{document}
\begin{abstract}
We establish the global isoperimetric inequality for convex geodesic tetrahedra in the unit three-sphere: at prescribed volume, surface area is minimized uniquely by the regular tetrahedron. The argument works equivalently at fixed surface area, where it maximizes volume. We first classify all interior constrained critical points by an analytic ordering lemma for the inverse Gram matrix of a spherical triangle. We then compactify the space of tetrahedra by their defining hemispheres and exclude every positive-volume degenerate maximizer. In the only nontrivial boundary stratum, the problem reduces to a spherical polygonal inequality and an explicit perturbation of a regular-square suspension.
\end{abstract}
\maketitle

\section{Introduction}
Let $\Sph^3\subset\R^4$ carry its unit-curvature round metric. A \emph{spherical tetrahedron} is the intersection with $\Sph^3$ of a full-dimensional simplicial convex cone generated by four linearly independent rays and contained in some open hemisphere. Its edges are the minimizing geodesic arcs between its vertices. We write $A(T)$ for the two-dimensional area of its four faces and $V(T)$ for its three-dimensional volume.

B\"ohm's work \cite{Bohm} established local isoperimetric extremality for regular spherical simplices. The distinction between local and global optimality is relevant to subsequent work on solid angles; see \cite{AkopyanKarasev}. Here we consider the global question in dimension three.

\begin{theorem}[Global isoperimetric inequality]\label{thm:main}
For every nondegenerate convex spherical tetrahedron $T\subset\Sph^3$, let $R$ be the regular tetrahedron with $V(R)=V(T)$. Then
\begin{equation}\label{eq:main}
 A(T)\ge A(R),
\end{equation}
with equality if and only if $T$ is congruent to $R$.
\end{theorem}

The proof is most naturally formulated in the dual, fixed-area form.
\begin{theorem}[Dual formulation]\label{thm:dual}
At each prescribed surface area $A_0\in(0,4\pi)$, the regular spherical tetrahedron is the unique volume maximizer among spherical tetrahedra of area $A_0$.
\end{theorem}

The main analytic step is the rigidity of interior critical points (Theorem~\ref{thm:critical}). The global step requires a compactification: a sequence of tetrahedra may develop antipodal vertices while its volume remains positive. Taking limits of the four supporting hemispheres, rather than of the vertex coordinates, retains the correct geometric limit.

\section{First variation and Gram coordinates}\label{sec:variation}

For a nondegenerate tetrahedron $T$, let $v_1,\ldots,v_4\in\Sph^3$ be its vertices, and let $n_1,\ldots,n_4\in\Sph^3$ be the unit inward normals to the respective opposite faces. Thus
\begin{equation}\label{eq:dual-normals}
 \ip{n_i}{v_j}=0\quad (i\ne j),\qquad \ip{n_i}{v_i}>0.
\end{equation}
The normals are linearly independent. There is a unique $q\in\R^4$ satisfying $\ip{n_i}{q}=1$ for each $i$. Put
\begin{equation}\label{eq:height}
 h_i=\ip{q}{v_i}\in\R.
\end{equation}
For $i<j$, let $\ell_{ij}\in(0,\pi)$ denote the length of the edge $v_iv_j$ and $\theta_{ij}\in(0,\pi)$ its interior dihedral angle.

\begin{lemma}[First variations]\label{lem:firstvariation}
With the above conventions,
\begin{equation}\label{eq:firstvariation}
 \dd V=\frac12\sum_{i<j}\ell_{ij}\dd\theta_{ij},
 \qquad
 \dd A=\sum_{i<j}(h_i+h_j)\tan\frac{\ell_{ij}}2\dd\theta_{ij}.
\end{equation}
\end{lemma}
\begin{proof}
The volume identity is Schlaefli's formula in curvature $+1$; see, for example, \cite[\S3.2]{Vinberg}.

For the area identity, consider the spherical polar tetrahedron $T^\circ$ with vertices $n_i$. Spherical angle-sum identities give
\begin{equation}\label{eq:polar-area}
 A(T)+A(T^\circ)=4\pi.
\end{equation}
For a spherical triangle with sides $a,b,c$, let $E(a,b,c)$ be its area and let $\Delta>0$ be the square root of its vertex Gram determinant. Differentiation of the spherical cosine rule yields
\begin{equation}\label{eq:triangle-area-derivative}
 \frac{\partial E}{\partial a}
 =\tan\frac a2\frac{1+\cos a-\cos b-\cos c}{\Delta}.
\end{equation}
Fix a dihedral angle $\theta_{ij}$, and denote the complementary vertex indices by $k,l$. The polar edge $n_kn_l$ has length $\pi-\theta_{ij}$ and belongs to two polar faces. Substituting their Gram data into \eqref{eq:triangle-area-derivative}, the sum of these two derivatives is
\[
 \ip{q}{\frac{v_i-\cos\ell_{ij}\,v_j}{\sin\ell_{ij}}
            +\frac{v_j-\cos\ell_{ij}\,v_i}{\sin\ell_{ij}}}
 =(h_i+h_j)\tan\frac{\ell_{ij}}2.
\]
(The two displayed unit vectors are the conormals dual to the two polar faces, obtained directly by projection to the plane $\operatorname{span}(v_i,v_j)$.) Since $\dd(\pi-\theta_{ij})=-\dd\theta_{ij}$, differentiation of \eqref{eq:polar-area} gives the second formula of \eqref{eq:firstvariation}.
\end{proof}

The off-diagonal entries of the normal Gram matrix are $\ip{n_i}{n_j}=-\cos\theta_{kl}$, where $\{i,j,k,l\}=\{1,2,3,4\}$. Positive definiteness is an open condition. Thus the six dihedral angles form local coordinates on the space of tetrahedra modulo isometries. The volume gradient in these coordinates never vanishes. The area gradient also never vanishes: if every $h_i+h_j=0$, then all $h_i=0$, and the basis $v_1,\ldots,v_4$ forces $q=0$, contradicting $\ip{q}{n_i}=1$.

\begin{proposition}[Critical-point equations]\label{prop:stationary}
At any interior critical point of surface area subject to fixed volume (or of volume subject to fixed surface area), there is a nonzero constant $c\in\R$ such that
\begin{equation}\label{eq:stationary}
 h_i+h_j=c\,f(\ell_{ij}),\qquad
 f(\ell)=\ell\cot\frac\ell2,\qquad 1\le i<j\le4.
\end{equation}
\end{proposition}
\begin{proof}
By the nonvanishing gradients and the Lagrange multiplier rule, the two differentials in \eqref{eq:firstvariation} are proportional. Cancelling the positive factors $\tan(\ell_{ij}/2)$ gives \eqref{eq:stationary}. If $c=0$, all $h_i+h_j$ vanish, a contradiction.
\end{proof}

\section{A strict ordering lemma for spherical triangles}\label{sec:triangle}

Introduce the functions $\Psi,\phi:(-1,1)\to\R$ by
\begin{equation}\label{eq:phi}
 \Psi(x)=\frac{\arccos x}{\sqrt{1-x^2}}
 =\int_0^\infty\frac{\dd s}{s^2+2xs+1},
 \qquad \phi(x)=(1+x)\Psi(x).
\end{equation}
They satisfy $\phi(\cos\ell)=f(\ell)$. Differentiation under the integral gives
\begin{equation}\label{eq:phi-concave}
 \phi'(x)=\int_0^\infty\frac{(s-1)^2}{(s^2+2xs+1)^2}\dd s>0,
 \qquad
 \phi''(x)=-4\int_0^\infty\frac{s(s-1)^2}{(s^2+2xs+1)^3}\dd s<0.
\end{equation}
Moreover, $\phi(-1+)=0$, $\phi(1-)=2$, and $\phi'(1-)=1/3$. Consequently, the inverse function
\begin{equation}\label{eq:g}
 g:(0,2)\longrightarrow(-1,1),\qquad g=\phi^{-1},
\end{equation}
is increasing and strictly convex. In particular, $\phi'(x)>1/3$ for $x<1$.

Let $\mathcal D\subset\R^3$ consist of all triples $t=(t_1,t_2,t_3)$ such that $t_i+t_j\in(0,2)$ for $i\ne j$ and the symmetric matrix
\begin{equation}\label{eq:K}
 K(t)=\begin{pmatrix}
 1&g(t_1+t_2)&g(t_1+t_3)\\
 g(t_1+t_2)&1&g(t_2+t_3)\\
 g(t_1+t_3)&g(t_2+t_3)&1
 \end{pmatrix}
\end{equation}
is positive definite. Define the real-valued function
\begin{equation}\label{eq:F}
 F:\mathcal D\to\R,\qquad F(t)=t^TK(t)^{-1}t.
\end{equation}

\begin{lemma}[Triangle ordering lemma]\label{lem:triangle}
The intersection of $\mathcal D$ with every line parallel to a coordinate axis is an interval. Furthermore:
\begin{enumerate}[label=\textup{(\roman*)},itemsep=2pt]
\item For $t\in\mathcal D$, if $t_i\ge t_j$ for some $j\ne i$, then $\partial_iF(t)<0$.
\item If $(b,a,a)\in\mathcal D$ with $b\le a$, then $\dfrac{\dd}{\dd b}F(b,a,a)<0$.
\end{enumerate}
\end{lemma}

\subsection{Inverse ordering}
Write $z=K(t)^{-1}t\in\R^3$. We first prove
\begin{equation}\label{eq:reverse}
 t_i>t_j\quad\Longrightarrow\quad z_i<z_j.
\end{equation}
Indeed, the matrix $K(t)$ is the Gram matrix of the vertices of a nondegenerate spherical triangle with side lengths $a,b,c\in(0,\pi)$, opposite the first, second, and third vertices. Set $X=\cos a$, $Y=\cos b$, $Z=\cos c$, and $\Delta^2=\det K(t)>0$. Then
\begin{equation}\label{eq:tsides}
 t_1=\tfrac12(f(b)+f(c)-f(a)),\quad
 t_2=\tfrac12(f(a)+f(c)-f(b)),\quad
 t_3=\tfrac12(f(a)+f(b)-f(c)).
\end{equation}
If $t_1>t_2$, then $a>b$. Put
\begin{equation}\label{eq:uvw}
 p=\frac{a+b}2,\quad q_0=\frac{a-b}2,\quad r=\frac c2,
 \qquad u=\cos p,\quad v=\cos q_0,\quad w=\cos r.
\end{equation}
The spherical triangle inequalities, including $a+b+c<2\pi$, imply
$0<q_0<r<\min(p,\pi-p)$, hence $|u|<w<v<1$.

Let $\eta(s)=(s^2+1)/(2s)\ge1$. Direct inversion of the three-by-three Gram matrix, followed by \eqref{eq:phi}, yields the exact identity
\begin{equation}\label{eq:kernel}
 z_1-z_2=-\frac{Y-X}{\Delta^2}\int_0^\infty
 \frac{\eta(s)(w-u)(v-w)(\eta(s)w-uv)}
 {s(\eta(s)+u)(\eta(s)+v)(\eta(s)+w)}\dd s.
\end{equation}
The integral is strictly positive: $w-u>0$, $v-w>0$, and $\eta(s)w-uv>0$. This proves \eqref{eq:reverse}. The identity is obtained by elementary rational simplification; the intermediate identities are recorded in Appendix~\ref{sec:algebra} for reproducibility.

\subsection{Positivity and derivatives}
We next show that every coordinate of $z$ is positive, with the uniform bound
\begin{equation}\label{eq:zlower}
 z_i\ge\frac13\qquad(i=1,2,3).
\end{equation}
First suppose $t_1>t_2>t_3$. By \eqref{eq:reverse}, $z_1<z_2<z_3$. Subtracting the second and third equations of $K(t)z=t$ gives
\begin{equation}\label{eq:rowdiff}
 t_2-t_3=(1-K_{23})(z_2-z_3)+(K_{12}-K_{13})z_1.
\end{equation}
The first term is strictly negative. Since $t_2-t_3=\phi(K_{12})-\phi(K_{13})$, we obtain
\[
 z_1>\frac{\phi(K_{12})-\phi(K_{13})}{K_{12}-K_{13}}>\frac13.
\]
The strict inequality holds by \eqref{eq:phi-concave}. Since $z_1$ is the smallest coordinate, all coordinates exceed $1/3$. Continuity gives \eqref{eq:zlower} without the distinctness assumption.

Differentiation of \eqref{eq:F}, using $\partial_iK_{ij}=g'(t_i+t_j)$ for $j\ne i$, gives
\begin{equation}\label{eq:dF}
 \partial_iF(t)=2z_i\left(1-\sum_{j\ne i}g'(t_i+t_j)z_j\right).
\end{equation}
If $t_i>t_j$ and $k$ is the remaining index, subtraction of rows $i$ and $j$ gives
\[
 t_i-t_j=(1-K_{ij})(z_i-z_j)+(K_{ik}-K_{jk})z_k
 <(K_{ik}-K_{jk})z_k.
\]
Since $g$ is convex,
\[
 g'(t_i+t_k)\ge\frac{K_{ik}-K_{jk}}{t_i-t_j},
\]
so $g'(t_i+t_k)z_k>1$. Formula \eqref{eq:dF}, together with \eqref{eq:zlower}, proves $\partial_iF<0$. When $t_i=t_j$, perturb $t_i$ slightly upward within the open domain and pass to the limit. Then $g'(t_i+t_k)z_k\ge1$, and the other positive summand in \eqref{eq:dF} still yields $\partial_iF<0$. This proves Lemma~\ref{lem:triangle}(i).

\subsection{The isosceles direction}
Consider $t=(b,a,a)$ with $b\le a$, and write $z=(U,W,W)$. Define $x=g(2a)$, $y=g(a+b)$, and $D=1+x-2y^2>0$. Thus $x\ge y$ and $x>2y^2-1$. Inverting $K$ gives
\begin{equation}\label{eq:UW}
 U=\frac{(1+x)b-2ya}{D},\qquad
 W=\frac{a-yb}{D}.
\end{equation}
Perturbing one of the two equal $a$-coordinates upward, and passing to the limit in the preceding row-difference argument, shows that
\begin{equation}\label{eq:Ulower}
 g'(a+b)U\ge1.
\end{equation}
We claim that $U<2W$. As $2a=\phi(x)$ and $a+b=\phi(y)$, \eqref{eq:UW} gives
\begin{equation}\label{eq:integral-iso}
 2D(U-2W)=2(1+x+2y)(1+y)\Psi(y)
 -(3+x+4y)(1+x)\Psi(x).
\end{equation}
Use the integral representation in \eqref{eq:phi}. With $h=\eta(s)\ge1$, the numerator of the combined integrand on the right of \eqref{eq:integral-iso}, over the positive denominator $(h+x)(h+y)$, equals
\begin{equation}\label{eq:P}
 (h-1)P(x,y)+(x^2-1)(1+y),
 \quad P(x,y)=-x^2-2xy-2x+4y^2+2y-1.
\end{equation}
Both terms are negative. Indeed, $x\ge y$ and $x>2y^2-1$ imply $\partial_xP=-2(x+y+1)<0$ throughout the domain. If $y\ge-\tfrac12$, then
$P(x,y)\le P(y,y)=y^2-1<0$; if $y<-\tfrac12$, then
\[
 P(x,y)<P(2y^2-1,y)=-4y(y-1)(y+1)^2<0.
\]
Also $x^2-1<0$ and $1+y>0$. Hence $U<2W$. Combining this with \eqref{eq:Ulower} gives $2g'(a+b)W>1$. Formula \eqref{eq:dF} therefore yields
\[
 \frac{\dd}{\dd b}F(b,a,a)=2U(1-2g'(a+b)W)<0,
\]
proving Lemma~\ref{lem:triangle}(ii).

\subsection{Coordinate intervals}
Let $L=f^{-1}:(0,2)\to(0,\pi)$. Since
\[
 f''(\ell)=\csc^2\frac\ell2\left(\frac{f(\ell)}2-1\right)<0,
\]
the inverse $L$ is decreasing and concave. Fix the second and third coordinates $y\ge z$, and vary the first coordinate $x$. The three side lengths of $K(x,y,z)$ are
\[
 a=L(y+z),\qquad b=L(x+z)\ge c=L(x+y).
\]
As $x$ increases, $b+c$ strictly decreases and $b-c$ is nondecreasing. The positive definiteness of the vertex Gram matrix is equivalent to the spherical triangle inequalities
\begin{equation}\label{eq:triangle-domain}
 b-c<a<b+c,\qquad a+b+c<2\pi,
\end{equation}
together with $a,b,c\in(0,\pi)$. Each inequality specifies an interval of admissible $x$-values. Intersecting these intervals with the two pair-sum conditions proves the interval assertion in Lemma~\ref{lem:triangle}.

\section{Rigidity of interior critical tetrahedra}\label{sec:critical}

\begin{theorem}\label{thm:critical}
Every nondegenerate constrained critical point of either $A$ at fixed $V$ or $V$ at fixed $A$ is a regular spherical tetrahedron.
\end{theorem}
\begin{proof}
Take the constant $c\ne0$ from \eqref{eq:stationary} and introduce four real parameters $t_i=h_i/c$. From \eqref{eq:stationary} and \eqref{eq:g},
\begin{equation}\label{eq:vertex-gram}
 \ip{v_i}{v_j}=g(t_i+t_j)\qquad(i\ne j).
\end{equation}
Thus the vertex Gram matrix of each face is $K$ evaluated at its three associated parameters.

For the face opposite $v_i$, the orthogonal projection of $q$ to its linear span is $q-n_i$. Its scalar products with the three face vertices equal the corresponding $h_j$. Consequently, the squared norm of $(q-n_i)/c$ is expressed by the quadratic form \eqref{eq:F}, giving
\begin{equation}\label{eq:fourF}
 F((t_j)_{j\ne i})=\frac{|q|^2-1}{c^2},
 \qquad i=1,\ldots,4.
\end{equation}

Relabel vertices so that $t_1\ge t_2\ge t_3\ge t_4$. First compare $F(t_1,t_3,t_4)$ and $F(t_2,t_3,t_4)$. By the coordinate-interval assertion in Lemma~\ref{lem:triangle}, the connecting interval is admissible. The function $F$ is strictly decreasing as the first coordinate varies there, since it is always at least $t_3$. Equality in \eqref{eq:fourF} implies $t_1=t_2$.

Similarly, compare $F(t_1,t_2,t_4)$ and $F(t_1,t_3,t_4)$. Throughout the relevant interval, the varying coordinate is at least $t_4$, so the lemma gives $t_2=t_3$. Write their common value as $a$ and put $b=t_4\le a$. Now \eqref{eq:fourF} reads $F(a,a,b)=F(a,a,a)$. The strict isosceles derivative from Lemma~\ref{lem:triangle}(ii), and the coordinate-interval property, force $b=a$.

Hence all four parameters $t_i$ coincide. Since $t_i+t_j=f(\ell_{ij})$, and $f$ is strictly decreasing, all six edge lengths are equal. Therefore $T$ is regular.
\end{proof}

\section{Compactness and boundary geometries}\label{sec:boundary}

Fix $A_0\in(0,4\pi)$. There is a regular tetrahedron with area $A_0$ (see Section~\ref{sec:completion}), so the supremum of the volumes of area-$A_0$ tetrahedra is positive.

Let $T_m$ be a maximizing sequence. Passing to a subsequence, the four inward unit normals converge to $n_1,\ldots,n_4\in\Sph^3$. Define the closed cone
\begin{equation}\label{eq:cone}
 C=\{x\in\R^4:\ip{n_i}{x}\ge0\ \text{for }i=1,\ldots,4\}.
\end{equation}
The intersection $C\cap\Sph^3$ has the limiting volume, because the indicators converge almost everywhere on $\Sph^3$ (the exceptional set is contained in four great spheres). In particular, $C$ has nonempty interior.

To obtain convergence of area as well, consider the convex bodies $K_m=C_m\cap B^4$ and $K=C\cap B^4$, where $B^4$ is the unit Euclidean ball. They converge in Hausdorff distance. Euclidean boundary measure in dimension three is continuous under this convergence; for instance, use the Cauchy projection formula for convex bodies \cite{Schneider}. Radial integration gives the exact identity
\begin{equation}\label{eq:radial-area}
 \mathcal H^3(\partial(C\cap B^4))
 =V(C\cap\Sph^3)+\frac13A(C\cap\Sph^3).
\end{equation}
The spherical boundary area of $C\cap\Sph^3$ is therefore the limit of $A(T_m)$, namely $A_0$. Thus the volume supremum is attained among these possibly degenerate cone sections.

Let $r$ be the dimension of the span of the four limiting normals. If $r=4$, then $C\cap\Sph^3$ is a genuine tetrahedron. If $r\le2$ and the volume is positive, its section is a hemisphere or a spherical wedge. In either case the boundary has area $4\pi$, contradicting $A_0<4\pi$. Hence only $r=3$ requires consideration.

In that case, after an orthogonal change of coordinates, $C=\R\times D$, where $D\subset\R^3$ is a full-dimensional pointed convex cone with at most four facets. Its spherical section $P=D\cap\Sph^2$ is a convex spherical polygon contained in an open hemisphere, with at most four sides. Parametrize the suspension by $(\cos t\,y,\sin t)$, for $y\in P$ and $-\pi/2\le t\le\pi/2$. Integrating the warped metric gives
\begin{equation}\label{eq:suspension}
 A(C\cap\Sph^3)=2\,\operatorname{per}(P),
 \qquad V(C\cap\Sph^3)=\frac\pi2\,\area(P).
\end{equation}

The spherical polygonal isoperimetric inequality, in the form allowing at most four sides \cite{DattaGupta}, implies that at fixed perimeter $\operatorname{per}(P)<2\pi$ its area is at most that of the regular spherical square with the same perimeter. Hence \emph{every rank-three boundary maximizer is dominated in volume by a regular-square suspension with exactly the same surface area}. It remains to show that such a suspension can always be improved by a nondegenerate tetrahedron.

\section{An explicit perturbation of the square suspension}\label{sec:perturb}

\begin{proposition}\label{prop:boundary}
For every $A_0\in(0,4\pi)$, the regular-square suspension of surface area $A_0$ admits a nondegenerate tetrahedral perturbation having exactly the same surface area and strictly greater volume.
\end{proposition}
\begin{proof}
Fix $a\in(0,1)$ and a small parameter $\delta\ge0$. Define $b=2a-1+\delta$ and the symmetric matrix
\begin{equation}\label{eq:gram-family}
 H(a,\delta)=\begin{pmatrix}
 1&a&b&a\\ a&1&a&b\\ b&a&1&a\\ a&b&a&1
 \end{pmatrix}.
\end{equation}
Its eigenvalues are $4a+\delta$, $\delta$, and $2(1-a)-\delta$ (with multiplicity two). For $\delta>0$ sufficiently small it is the Gram matrix of four independent inward face normals and hence defines a nondegenerate spherical tetrahedron. When $\delta=0$ it defines a regular-square suspension.

Let $A(a,\delta)$ and $V(a,\delta)$ be its spherical surface area and volume. The square cross-section has side length $2\arctan\sqrt a$ and each interior angle equal to $\pi/2+\arcsin a$. Equation \eqref{eq:suspension} gives
\begin{equation}\label{eq:boundary-base}
 A(a,0)=16\arctan\sqrt a,
 \qquad V(a,0)=2\pi\arcsin a.
\end{equation}
As $a$ runs through $(0,1)$, these surface areas fill $(0,4\pi)$.

Every polar face has side cosines $a,a,b$. Its Gram determinant is
$\Delta(a,\delta)^2=(1-b)(1+b-2a^2)$. The spherical triangle area is
\[
 E(a,\delta)=2\arctan\frac{\Delta(a,\delta)}{1+2a+b}.
\]
By polarity, $A=4\pi-4E$. Direct differentiation at fixed $a$ gives
\begin{equation}\label{eq:Aexact}
 \frac{\partial A}{\partial\delta}
 =\frac{4\delta}{(1+b)\Delta(a,\delta)}.
\end{equation}
Thus the area change has no linear term. Taylor expansion at $\delta=0$ yields
\begin{equation}\label{eq:Aexpansion}
 A(a,\delta)=A(a,0)
 +\frac{\delta^2}{2a^{3/2}(1-a)}+O_a(\delta^3).
\end{equation}

At fixed $a$, the two dihedral angles corresponding to the entries $b$ vary, and the other four remain fixed. By Schlaefli's formula,
\begin{equation}\label{eq:Vexact}
 \frac{\partial V}{\partial\delta}
 =\frac{\ell_b(a,\delta)}{\sqrt{1-b^2}},
\end{equation}
where $\ell_b$ is the common length of the two original edges dual to these entries. To compute it, put
\[
 m=2(1-a)-\delta,
 \qquad S=\frac1{4a+\delta}+\frac1\delta+\frac2m.
\]
The normal Gram matrix is diagonalized by the four sign vectors with entries $\pm1$. Normalizing $H^{-1}$ to unit diagonal gives the vertex Gram matrix, hence
\[
 \cos\ell_b=1-\frac4{mS},\qquad
 \ell_b=2\sqrt{\frac\delta{1-a}}+O_a(\delta^{3/2}).
\]
Since $\sqrt{1-b^2}=2\sqrt{a(1-a)}+O_a(\delta)$, integration of \eqref{eq:Vexact} yields
\begin{equation}\label{eq:Vexpansion}
 V(a,\delta)=V(a,0)
 +\frac{2\delta^{3/2}}{3\sqrt a(1-a)}
 +O_a(\delta^{5/2}).
\end{equation}

Fix $a_0\in(0,1)$. As
$\partial_aA(a,0)=8/[\sqrt a(1+a)]>0$, the implicit function theorem applied to the smooth area function gives a smooth parameter $a(\delta)=a_0+O(\delta^2)$ satisfying
\[
 A(a(\delta),\delta)=A(a_0,0)
\]
for all sufficiently small $\delta\ge0$. From \eqref{eq:boundary-base} and \eqref{eq:Vexpansion},
\begin{equation}\label{eq:gain}
 V(a(\delta),\delta)-V(a_0,0)
 =\frac{2\delta^{3/2}}{3\sqrt{a_0}(1-a_0)}
 +O_{a_0}(\delta^2)>0.
\end{equation}
This is the required fixed-area improvement.
\end{proof}

\begin{corollary}\label{cor:interior}
For every $A_0\in(0,4\pi)$, the maximal volume of a spherical tetrahedron of surface area $A_0$ is attained by a nondegenerate tetrahedron.
\end{corollary}
\begin{proof}
The compactification in Section~\ref{sec:boundary} supplies a maximizer. The rank-one and rank-two strata have area $4\pi$ and cannot occur. Every rank-three stratum is bounded above by a regular-square suspension of the same area, which Proposition~\ref{prop:boundary} strictly improves. Therefore the maximizer is interior.
\end{proof}

\section{Completion of the proof}\label{sec:completion}

\begin{proof}[Proof of Theorems~\ref{thm:dual} and \ref{thm:main}]
Fix $A_0\in(0,4\pi)$. By Corollary~\ref{cor:interior}, the volume supremum at fixed area is attained by a nondegenerate tetrahedron. This tetrahedron is an interior constrained critical point; Theorem~\ref{thm:critical} forces it to be regular. All regular tetrahedra with a given surface area are congruent, as seen from the following formula.

For a regular spherical tetrahedron $R_\ell$ with common edge length $\ell\in(0,\arccos(-1/3))$, the spherical cosine rule gives the angle of each equilateral face as $\arccos(\cos\ell/(1+\cos\ell))$. Thus
\begin{equation}\label{eq:regular-area}
 A(R_\ell)=12\arccos\!\left(\frac{\cos\ell}{1+\cos\ell}\right)-4\pi.
\end{equation}
This is strictly increasing from $0$ to $4\pi$. The internal dihedral angle $\theta(\ell)$ satisfies
\begin{equation}\label{eq:regular-dihedral}
 \cos\theta(\ell)=\frac{\cos\ell}{1+2\cos\ell}.
\end{equation}
It too increases strictly with $\ell$. By Schlaefli's formula,
\[
 \frac{\dd}{\dd\ell}V(R_\ell)=3\ell\,\theta'(\ell)>0.
\]
As $\ell\to0$ the tetrahedron contracts and its volume tends to $0$; as $\ell\to\arccos(-1/3)$ its four normals coalesce and the tetrahedron tends to a hemisphere of volume $\pi^2$. Hence the regular family provides a strictly increasing bijection from edge length onto the volume interval $(0,\pi^2)$, as well as onto the area interval $(0,4\pi)$.

We have shown that for every area $A_0$, the regular tetrahedron uniquely maximizes volume. The strict monotonicity of both quantities along the regular family converts this inequality into \eqref{eq:main}, with the stated equality condition.
\end{proof}

\appendix
\section{Two algebraic identities behind the triangle lemma}\label{sec:algebra}

This appendix records the rational simplifications used in Section~\ref{sec:triangle}. They are identities, independent of any inequality assumption.

For $X,Y,Z\in(-1,1)$, define $\phi_X=\phi(X)$, $\phi_Y=\phi(Y)$, $\phi_Z=\phi(Z)$ and $m=(\phi_Y-\phi_X)/(Y-X)$. The inverse of
$K=\begin{psmallmatrix}1&Z&Y\\Z&1&X\\Y&X&1\end{psmallmatrix}$
gives
\begin{align}\label{eq:appendix-J}
 2(\det K)(z_1-z_2)&=(Y-X)J,\\
 J&=(1+X+Y+Z)\phi_Z
 -(1+Z)(\phi_X+\phi_Y-2m)-(X+Y)^2m.\notag
\end{align}
Under the substitutions \eqref{eq:uvw}, the half-angle formulas give
\begin{equation}\label{eq:appendix-J2}
 \frac J4=w(w^2+uv)\Psi(w)
 +\frac{u^2(w^2-v^2)\Psi(u)+v^2(u^2-w^2)\Psi(v)}{v-u}.
\end{equation}
For any $h\in\R$ such that the denominators are nonzero,
\begin{align*}
 &\frac{w(w^2+uv)}{h+w}
 +\frac{u^2(w^2-v^2)/(h+u)+v^2(u^2-w^2)/(h+v)}{v-u}\\
 &\hspace{22mm}=-\frac{h(u-w)(v-w)(uv-hw)}{(h+u)(h+v)(h+w)}.
\end{align*}
Using $\Psi(t)=\int_0^\infty[2s(\eta(s)+t)]^{-1}\dd s$ in \eqref{eq:appendix-J2} yields \eqref{eq:kernel}.

For the isosceles argument, set $P(x,y)$ as in \eqref{eq:P}. The second identity is
\begin{align*}
 &\frac{2(1+x+2y)(1+y)}{h+y}
 -\frac{(3+x+4y)(1+x)}{h+x}\\
 &\hspace{16mm}
 =\frac{(h-1)P(x,y)+(x^2-1)(1+y)}{(h+x)(h+y)}.
\end{align*}
Both identities can be checked by clearing denominators.

\begin{thebibliography}{9}
\bibitem{AkopyanKarasev}
A.~Akopyan and R.~Karasev,
\emph{Bounding minimal solid angles of polytopes},
arXiv:1505.05263 (2015), revised 2016.
\url{https://arxiv.org/abs/1505.05263}.

\bibitem{Bohm}
J.~B\"ohm,
\emph{Zum isoperimetrischen Problem bei Simplexen in $r$-dimensionalen R\"aumen konstanter Kr\"ummung},
Studia Sci. Math. Hungar. \textbf{21} (1986), no.~1--2, 267--281.

\bibitem{DattaGupta}
B.~Datta and S.~Gupta,
\emph{Isoperimetric inequality for non-Euclidean polygons},
arXiv:2409.06529 (2024).
\url{https://arxiv.org/abs/2409.06529}.

\bibitem{Schneider}
R.~Schneider,
\emph{Convex Bodies: The Brunn--Minkowski Theory},
2nd expanded ed., Encyclopedia of Mathematics and its Applications 151,
Cambridge University Press, 2014.

\bibitem{Vinberg}
\`E.~B.~Vinberg,
\emph{Volumes of non-Euclidean polyhedra},
Russian Math. Surveys \textbf{48} (1993), no.~2, 15--45.
\url{https://doi.org/10.1070/RM1993v048n02ABEH001011}.
\end{thebibliography}
\end{document}
